% !TEX root = TFM-ML-01-memoria-referencia.tex
\section{Investigación de usuario}\label{sec:investigacion-usuario}

\subsection{Diseño y metodología de la encuesta primaria}
\label{sec:metodologia-encuesta}

La investigación primaria de este trabajo consiste en
una encuesta dirigida a trabajadores sanitarios de turno
nocturno en hospitales de la provincia de A~Coruña,
diseñada y administrada por el autor en el marco del
presente TFM. A continuación se describe su ficha técnica
y las decisiones metodológicas que la sustentan.

\subsubsection{Ficha técnica}
\label{sec:ficha-tecnica}

\begin{itemize}
  \item \textbf{Universo:} Personal sanitario que trabaja
    o ha trabajado en turno de noche en hospitales de la
    provincia de A~Coruña.
  \item \textbf{Muestra válida:} 167~respuestas.
  \item \textbf{Período de recogida:} Marzo de 2026.
  \item \textbf{Instrumento:} Cuestionario estructurado
    de 25~ítems, organizado en seis bloques temáticos:
    perfil y demografía, patrón de desplazamiento, vacío
    de transporte público, carga del desplazamiento,
    impacto laboral, y fatiga y seguridad.
  \item \textbf{Distribución:} Canal digital. Difusión
    a través de redes profesionales del sector sanitario.
  \item \textbf{Análisis cuantitativo:} Microsoft Excel
    y JASP. Estadística descriptiva con frecuencias
    absolutas, porcentajes e índices de gravedad por ítem.
  \item \textbf{Análisis cualitativo:} 46~respuestas
    abiertas codificadas en 8~códigos temáticos mediante
    análisis de contenido inductivo.
\end{itemize}

\subsubsection{Limitaciones metodológicas}
\label{sec:limitaciones-encuesta}

La muestra es de conveniencia y no probabilística.
Los resultados no son extrapolables estadísticamente al
conjunto del personal sanitario nocturno de la provincia,
pero son suficientes para identificar patrones y
fundamentar criterios de diseño. La distribución digital
introduce un sesgo de autoselección que no puede
cuantificarse con los datos disponibles
\cite{LouridoRegueira2026}.

\subsection{Perfil del usuario primario: personal sanitario nocturno}
\label{sec:perfil-usuario}

El perfil que se presenta a continuación se construye
íntegramente desde los datos de la encuesta propia.
La literatura europea se utiliza como marco de
contextualización, no como fuente del perfil.

\subsubsection{Categoría profesional}
\label{sec:categoria-profesional}

El colectivo mayoritario es el de enfermería, que
representa el 56{,}9\,\% de la muestra (95 de~167).
El segundo grupo es el de auxiliares de enfermería
(TCAE), con el 21{,}0\,\% (35 de~167). Ambos colectivos
suman el 77{,}9\,\% del total. Los médicos residentes
o adjuntos representan el 15{,}0\,\% (25 de~167),
y los celadores el 4{,}8\,\% (8 de~167)
\cite{LouridoRegueira2026}.

\subsubsection{Régimen de turno y noches al mes}
\label{sec:regimen-turno}

El 81{,}4\,\% de los encuestados (136 de~167) trabaja
en régimen de turno rotatorio con noches. El 49{,}1\,\%
(82 de~167) trabaja entre 5 y~8 noches al mes, y el
47{,}9\,\% (80 de~167) entre 1 y~4 noches. Solo el
2{,}9\,\% supera las 8~noches mensuales
\cite{LouridoRegueira2026}.

\subsubsection{Zona de residencia}
\label{sec:zona-residencia}

El 38{,}9\,\% de los encuestados (65 de~167) reside en
un municipio distinto al del hospital. El 32{,}9\,\%
(55 de~167) reside en el centro urbano de la misma
ciudad. El 17{,}9\,\% (30 de~167) reside en zona rural,
y el 10{,}2\,\% (17 de~167) en un barrio periférico de
la misma ciudad. El 49{,}1\,\% reside, por tanto, fuera
del entorno urbano inmediato del hospital
\cite{LouridoRegueira2026}.

\subsubsection{Género y edad}
\label{sec:genero-edad}

El 91{,}0\,\% de la muestra se identifica como mujer
(152 de~167). El 7{,}8\,\% se identifica como hombre
(13 de~167). El rango de edad predominante es el de
46--55~años, con el 34{,}7\,\% (58 de~167), seguido
de los rangos 26--35~años (21{,}0\,\%) y 36--45~años
(20{,}4\,\%) \cite{LouridoRegueira2026}.

\subsubsection{Hospital de referencia}
\label{sec:hospital-referencia}

El 44{,}3\,\% de los encuestados (74 de~167) trabaja
en el CHUF --- Complexo Hospitalario Universitario de
Ferrol. El 27{,}5\,\% (46 de~167) trabaja en el CHUAC
de A~Coruña, y el 20{,}4\,\% (34 de~167) en el CHUS
de Santiago de Compostela. El CHUF es el hospital con
mayor representación en la muestra y el caso de estudio
principal del trabajo \cite{LouridoRegueira2026}.

\subsection{Hallazgos cuantitativos: barreras de movilidad documentadas}
\label{sec:hallazgos-cuantitativos}

Los seis bloques del cuestionario arrojan los siguientes
resultados sobre las barreras de movilidad nocturna
del personal sanitario encuestado.

\subsubsection{Patrón de desplazamiento actual}
\label{sec:patron-desplazamiento}

El vehículo privado es el modo de transporte principal
tanto para ir al hospital como para volver. La dependencia
del vehículo propio en la franja nocturna refleja la
ausencia de alternativas, no necesariamente una preferencia
modal \cite{LouridoRegueira2026}.

\subsubsection{Vacío de transporte público}
\label{sec:vacio-tp}

El 52{,}1\,\% de los encuestados (87 de~167) carece de
cobertura fiable de transporte público en la franja
de su turno nocturno: 39 declaran no tener servicio,
y 48 declaran no saber si existe. El 47{,}3\,\%
(79 de~167) no encuentra transporte público disponible
\emph{varias veces a la semana}. El 30{,}5\,\%
(51 de~167) identifica la falta de transporte público
como el cambio prioritario único en su desplazamiento
\cite{LouridoRegueira2026}.

\subsubsection{Carga del desplazamiento}
\label{sec:carga-desplazamiento}

Las dificultades de aparcamiento afectan al 61{,}1\,\%
de los encuestados (102 de~167) con una frecuencia de
\emph{varias veces a la semana}. El índice de gravedad
de este ítem es de 60{,}9 sobre~100. El 26{,}9\,\%
(45 de~167) señala el aparcamiento garantizado como
el cambio prioritario único \cite{LouridoRegueira2026}.

\subsubsection{Impacto laboral}
\label{sec:impacto-laboral}

Los problemas de transporte nocturno tienen consecuencias
laborales concretas y cuantificables. El 6{,}6\,\% de los
encuestados (11 de~167) ha rechazado o evitado al menos
un turno de noche por dificultades de desplazamiento.
El 22{,}8\,\% (38 de~167) ha llegado tarde al turno de
noche por la misma causa al menos una vez. En total,
el 25{,}7\,\% de la muestra (43 de~167) ha experimentado
al menos uno de estos dos impactos laborales directos
\cite{LouridoRegueira2026}.

\subsubsection{Fatiga al volante}
\label{sec:fatiga-volante}

El 52{,}7\,\% de los encuestados (88 de~167) declara
conducir con sueño o fatiga \emph{varias veces a la
semana}. El índice de gravedad de este ítem es de
60{,}2 sobre~100, el más alto de los ítems de problema
del cuestionario \cite{LouridoRegueira2026}.

\subsection{Hallazgos cualitativos: análisis de respuestas abiertas}
\label{sec:hallazgos-cualitativos}

El cuestionario incluía una pregunta abierta opcional al
final. 46 de los 167 encuestados respondieron, lo que
representa una tasa de respuesta cualitativa del 27{,}5\,\%.
Las respuestas fueron codificadas por el autor en
8~códigos temáticos mediante análisis de contenido
inductivo \cite{LouridoRegueira2026}.

\subsubsection{Patrones de co-ocurrencia}
\label{sec:patrones-coocurrencia}

21 de las 46~respuestas abiertas (45{,}7\,\%) combinan
dos o más problemas simultáneos. Este patrón de
co-ocurrencia indica que las barreras de movilidad
nocturna no son independientes: la falta de transporte
público, la fatiga y el coste del aparcamiento se
experimentan de forma simultánea y se refuerzan
mutuamente \cite{LouridoRegueira2026}.

\subsubsection{Propuesta espontánea de servicio}
\label{sec:propuesta-espontanea}

Una respuesta, codificada como T8, formula espontáneamente
la demanda de un servicio de shuttle coordinado por la
institución sanitaria, sin que el cuestionario lo indujera
ni sugiriera:

\begin{quote}
«Estaría bien que a los trabajadores de hospitales se nos
pusiera un transporte para llevarnos al hospital y
recogernos en un punto y de ahí para casa.»
\hfill\textit{\cite{LouridoRegueira2026}}
\end{quote}

Esta respuesta es significativa no por su frecuencia ---
aparece una sola vez --- sino porque formula de forma
autónoma un modelo de servicio que coincide con la
hipótesis de trabajo de este proyecto. Su valor radica
en que no es inducida \cite{LouridoRegueira2026}.

\subsection{Perspectiva de género: brecha en el riesgo de fatiga al volante}
\label{sec:perspectiva-genero}

El cruce de los datos de fatiga al volante por género
revela una brecha estadísticamente relevante que tiene
implicaciones directas para el diseño del servicio.

El 90{,}8\,\% de las mujeres de la muestra (138 de~152)
presentan riesgo de fatiga al volante con una frecuencia
igual o superior a \emph{varias veces al mes}. En el
caso de los hombres, ese porcentaje es del 53{,}8\,\%
(7 de~13) \cite{LouridoRegueira2026}.

La muestra es 91{,}0\,\% femenina (152 de~167). Esto
significa que el colectivo que más necesita una alternativa
al vehículo privado nocturno es, de forma abrumadora,
el mismo que compone la mayor parte de la plantilla
sanitaria nocturna. La perspectiva de género no es un
eje añadido al análisis: emerge directamente de los datos.

Un servicio de movilidad nocturna para este colectivo
debe contemplar la seguridad percibida como criterio
de diseño de primer orden, dado que el 91\,\% de
sus usuarias potenciales son mujeres que viajan solas
en franja nocturna \cite{LouridoRegueira2026}.

\subsection{Benchmarking de pilotos de vehículos autónomos compartidos}\label{sec:benchmarking-av}
\begin{pendiente}
%% INSERTAR: benchmarking-pilotos-av
%% FIGURA: fig:benchmarking-av
\end{pendiente}

\subsection{Benchmarking de servicios nocturnos en ciudades comparables}\label{sec:benchmarking-ciudades}
\begin{pendiente}
%% INSERTAR: benchmarking-ciudades-medias
\end{pendiente}
